\documentclass[11pt]{article}

\usepackage[margin=1in]{geometry}
\usepackage{amsmath,amssymb}
\usepackage{booktabs}
\usepackage{enumitem}
\usepackage{hyperref}
\usepackage{microtype}
\usepackage{xcolor}
\usepackage{listings}

\hypersetup{
  colorlinks=true,
  linkcolor=blue!70!black,
  citecolor=blue!70!black,
  urlcolor=blue!70!black
}

\lstset{
  basicstyle=\ttfamily\small,
  frame=single,
  columns=fullflexible,
  keepspaces=true,
  breaklines=true,
  showstringspaces=false,
  backgroundcolor=\color{black!2},
  rulecolor=\color{black!30},
  xleftmargin=2em,
  framexleftmargin=1.5em,
  numbers=left,
  numberstyle=\tiny\color{gray},
  numbersep=8pt,
  tabsize=2
}

% MeTTa listing language
\lstdefinelanguage{MeTTa}{
  morekeywords={if,let,let*,case,match,import!,bind!,assertEqual,
    assertEqualToResult,progn,prog1,collapse,println!},
  sensitive=true,
  morecomment=[l]{;;;},
  morecomment=[l]{;;},
  morestring=[b]",
  literate={!}{{\textbf{!}}}1
}

% Prolog listing language
\lstdefinelanguage{SWIProlog}{
  morekeywords={module,use_module,is,not,nb_setval,nb_getval,
    throw,catch,forall,findall,maplist,foldl,assoc_to_keys,
    put_assoc,get_assoc,ord_add_element,ord_subtract},
  sensitive=true,
  morecomment=[l]{\%},
  morecomment=[n]{/*}{*/},
  morestring=[b]",
  morestring=[b]'
}

\title{Metamath Proof Verification in MeTTa:\\
From Pure Prolog to Opaque and Space-Backed PeTTa Verifiers\\[6pt]
\large (AIrxiv Working Paper, 2026-02-25)}
\author{
  Zarathustra Goertzel\thanks{Human supervisor and original mmverify author.}
  \and
  Oru\v{z}i\thanks{AI-assisted collaboration
    (Claude Code, Codex; under human supervision).}
}
\date{February 2026}

\begin{document}
\maketitle

\begin{abstract}
We present \textbf{pverify}, a family of Metamath proof verifiers
developed in a staged methodology:
\textbf{(0)} \texttt{mmverify} (original PeTTa verifier),
\textbf{(1)} \texttt{pverify\_hybrid}
(HE MeTTa non-parsing stage based on the
\texttt{mmverify.py} parsing workflow),
\textbf{(2)} \texttt{pverify.pl} (pure Prolog baseline),
\textbf{(3)} \texttt{pverify\_op}
(verification semantics in PeTTa/MeTTa with opaque Prolog data
structures through FFI),
and \textbf{(4)} \texttt{pverify\_op\_space}
(space-backed storage using specialized
\texttt{MapSpace}/\texttt{SetSpace} backends).
All active implementations pass the same 141-test Metamath
conformance suite.

On \texttt{set.mm} (519K lines, 136{,}289 statements), controlled
solo benchmarks on the same machine yield:
5m\,58s for pure Prolog,
28m\,31s for \texttt{pverify\_op},
and 28m\,48s for \texttt{pverify\_op\_space}.
Thus the space-backed verifier is within 0.95\% wall-time of the
opaque-state verifier while preserving identical behavior.

\paragraph{Measurement provenance note.}
Performance/count claims in this draft are snapshot measurements. Primary artifacts:
\texttt{AIrxiv/airxiv\_2026-02-21\_pverify\_metamath\_in\_petta.log},
\texttt{AIrxiv/airxiv\_2026-02-21\_pverify\_metamath\_in\_petta.out},
\texttt{AIrxiv/airxiv\_2026-02-21\_pverify\_metamath\_in\_petta.fdb\_latexmk},
plus regression harnesses under
\texttt{hyperon/PeTTa/tests/test\_all\_mmverify.sh} and
\texttt{hyperon/PeTTa/tests/test\_all\_examples\_modified.sh}.

Methodologically, we make explicit what is and is not in MeTTa.
In \texttt{pverify\_op} and \texttt{pverify\_op\_space},
statement processing, RPN proof checking, substitution, and DV checks
are implemented in \texttt{.metta}; Prolog is used for parsing and
indexed primitives.
The specialized spaces are bonafide PeTTa spaces:
they support normal space operations
(\texttt{add-atom}, \texttt{remove-atom}, \texttt{match},
\texttt{get-atoms}) while exposing optimized indexed accessors for
verifier workloads.
The main engineering conclusion is that FFI granularity, not
algorithmic mismatch, remains the dominant optimization axis.
\end{abstract}

%% ================================================================
\section{Introduction}
\label{sec:intro}

Metamath~\cite{megill2019} is a minimalist proof language whose
simplicity---constants, variables, substitution, disjoint-variable
constraints, and a stack-based proof checker---makes it an
attractive target for implementation in new programming languages.

MeTTa~\cite{goertzel2023metta} is the primary language of the
Hyperon AGI framework.  It features non-deterministic evaluation,
pattern matching, and a space-based knowledge representation.
PeTTa~\cite{petta} is a Prolog-based implementation of MeTTa
that compiles MeTTa programs to SWI-Prolog, providing access to
Prolog's mature ecosystem through a foreign function interface (FFI).

This paper describes \textbf{pverify}, a family of Metamath proof
verifiers developed through human--AI collaboration.
The project evolved through several phases:

\begin{enumerate}[leftmargin=2em]
  \item \textbf{mmverify} (Goertzel): The original PeTTa Metamath
    verifier, hand-written, establishing the core verification
    algorithm in MeTTa with PeTTa-native data structures.
  \item \textbf{pverify\_hybrid} (AI agents): An HE MeTTa non-parsing
    precursor stage derived from the
    \texttt{mmverify.py} parsing approach.
  \item \textbf{pverify.pl} (AI, Claude Code): A pure Prolog
    reimplementation sharing the \texttt{mm\_primitives.pl} parser.
  \item \textbf{pverify\_op} (AI, Claude Code): The ``opaque data
    structures'' architecture---all verification logic in PeTTa,
    with Prolog providing only $O(\log n)$ AVL trees and ordsets
    through FFI.
  \item \textbf{pverify\_op\_stream} (AI, Claude Code + Codex):
    A streaming variant of pverify\_op that processes statements
    one-at-a-time from a Prolog-resident token stream, avoiding
    the need to hold all parsed statements in PeTTa memory
    simultaneously.
  \item \textbf{pverify\_op\_space} (AI, Claude Code + Codex):
    A space-backed variant of pverify\_op that replaces opaque
    Prolog assocs/ordsets with PeTTa-native \texttt{MapSpace} and
    \texttt{SetSpace} backends, achieving full parity while making
    the verifier's state backend-flexible.
\end{enumerate}

All current pverify variants share a common Prolog parsing layer
(\texttt{mm\_primitives.pl}, 1{,}224 lines) and are tested against
the same 141-test conformance suite covering both positive
(correct proofs accepted) and negative (invalid inputs rejected)
cases.

\subsection{Claims and Scope}
\label{sec:claims}

\textbf{What this paper claims.}
\begin{itemize}[leftmargin=2em]
  \item Four current \texttt{pverify} implementations pass the same
    141-test conformance suite.
  \item Controlled \texttt{set.mm} benchmarks are reported for the
    main methodology chain:
    \texttt{pverify.pl} $\rightarrow$
    \texttt{pverify\_op} $\rightarrow$
    \texttt{pverify\_op\_space}.
  \item \texttt{mmverify} and \texttt{pverify\_hybrid} are included
    as explicit historical/motivational precursors to this chain.
  \item \texttt{pverify\_op\_space} achieves both conformance parity
    and near wall-time parity with \texttt{pverify\_op}
    (0.95\% delta on \texttt{set.mm} in current measurements).
\end{itemize}

\textbf{What this paper does not claim.}
\begin{itemize}[leftmargin=2em]
  \item It is not a mechanized soundness/completeness proof of a
    verifier; this is addressed in separate Lean work.
  \item It does not claim per-error-code completeness (``first
    violated rule implies specific parser code'').
\end{itemize}

%% ================================================================
\section{Architecture}
\label{sec:arch}

\subsection{Methodology Chain and Invariants}
\label{sec:method-chain}

The experimental design follows one controlled chain:
\begin{enumerate}[leftmargin=2em]
  \item \textbf{Original PeTTa verifier}
    (\texttt{mmverify}): establishes the direct MeTTa implementation
    style and motivates optimization work.
  \item \textbf{HE MeTTa precursor}
    (\texttt{pverify\_hybrid}): HE MeTTa non-parsing verifier stage
    based on \texttt{mmverify.py} parser conventions.
  \item \textbf{Pure Prolog baseline}
    (\texttt{pverify.pl}): establishes a native lower-bound
    implementation over the same parser.
  \item \textbf{Opaque PeTTa implementation}
    (\texttt{pverify\_op}): moves verification semantics into MeTTa,
    with Prolog used only for parsing and indexed primitives.
  \item \textbf{Space-backed PeTTa implementation}
    (\texttt{pverify\_op\_space}): changes storage backend to
    specialized spaces while preserving verifier semantics.
\end{enumerate}

To keep this comparison honest, each stage preserves:
\begin{itemize}[leftmargin=2em]
  \item the same parser (\texttt{mm\_primitives.pl}),
  \item the same acceptance/rejection contract
    (141-test conformance suite),
  \item the same benchmark target (\texttt{set.mm}).
\end{itemize}

\subsection{Naive MeTTa-Style Baseline (Motivation)}
\label{sec:naive-metta}

A natural first implementation in MeTTa uses plain lists for maps and
sets with linear traversal:
\begin{lstlisting}[language=MeTTa,caption={Naive MeTTa-style lookup (linear scan).}]
;;; Labels as list of (pair Key Value)
(= (lookup-label $k $labels)
   (if (== (size-atom $labels) 0)
       not_found
       (let (($h $t) (decons-atom $labels))
         (let (($k0 $v0) $h)
           (if (== $k $k0) $v0 (lookup-label $k $t))))))
\end{lstlisting}

This style is simple and idiomatic, but it scales poorly when used in
hot loops (proof-step dispatch, substitution lookups, DV checks). The
methodology chain above is therefore a sequence of controlled moves
from naive linear structures to indexed backends while preserving
MeTTa-level verifier semantics.

\subsection{Shared Prolog Layer}

All verifiers in the pverify family share
\texttt{mm\_primitives.pl}, which provides:
\begin{itemize}[leftmargin=2em]
  \item Tokenization of \texttt{.mm} files (handling comments,
    includes, whitespace)
  \item Recursive-descent parsing into compound terms
    (\texttt{c([\ldots])}, \texttt{v([\ldots])}, etc.)
  \item File-wide semantic validation (duplicate labels, frame
    balance, typecode consistency, \texttt{\$d} variable checks,
    scoping rules)
  \item Streaming parser (\texttt{next\_statement/2}) that yields
    one statement at a time from a token list
\end{itemize}

\subsection{pverify\_op: Batch Verification in PeTTa}

The key architectural decision in \texttt{pverify\_op} is a clean
separation of concerns:

\begin{itemize}[leftmargin=2em]
  \item \textbf{Prolog (pverify\_ds.pl, 227 lines):} Provides
    data structure operations (AVL tree get/put, ordset
    add/member/union/subtract, compound term manipulation) as FFI
    exports.  These are \emph{opaque} to PeTTa---MeTTa code never
    inspects the internal structure of an AVL tree or ordset.
  \item \textbf{PeTTa (pverify\_op.metta, 693 lines):} Contains
    \emph{all} verification logic: statement processing
    (\texttt{\$c}, \texttt{\$v}, \texttt{\$f}, \texttt{\$e},
    \texttt{\$d}, \texttt{\$a}, \texttt{\$p}, scoping), the
    RPN proof stack machine (\texttt{treat-step}), substitution
    application (\texttt{apply-subst}), and
    disjoint-variable checking (\texttt{check-dvs}).
\end{itemize}

Verification state is threaded functionally as an immutable tuple:

\begin{lstlisting}[language=MeTTa,caption={State representation.}]
;;; State = (state Constants FrameStack Labels)
;;;   Constants:  opaque Prolog ordset
;;;   FrameStack: MeTTa list of Frame tuples
;;;   Labels:     opaque Prolog assoc (label -> entry)
\end{lstlisting}

Each statement processor receives the current state and returns
an updated state (or an error), with no mutable global variables.

\subsection{pverify\_op\_stream: Streaming Verification}

The streaming variant replaces batch parsing with incremental
statement delivery.  Three technical challenges arose:

\paragraph{Token list FFI boundary.}
Metamath token lists contain literal \texttt{(} and \texttt{)}
atoms.  PeTTa's FFI layer interprets parentheses as MeTTa grouping,
destructively altering the token structure.  The solution is to
keep the token stream entirely in Prolog memory using
\texttt{nb\_setval}/\texttt{nb\_getval} (SWI-Prolog non-backtrackable
globals), passing only parsed statement compounds across the FFI
boundary.

\begin{lstlisting}[language=SWIProlog,caption={Streaming helpers in pverify\_ds.pl.}]
ds_stream_init(Filename, ok) :-
    ( parse_mm_file_compounds(Filename, _) -> true
    ; throw(mm_error(parse_failed(Filename)))
    ),
    tokenize_mm_file(Filename, Tokens),
    nb_setval(mm_token_stream, Tokens).

ds_stream_next(Result) :-
    nb_getval(mm_token_stream, Tokens),
    next_statement(Tokens, NSResult),
    ( NSResult = [end]    -> Result = end
    ; NSResult = [error]  -> Result = error
    ; NSResult = [Stmt, Rest]
      -> nb_setval(mm_token_stream, Rest),
         Result = ns(Stmt)
    ).
\end{lstlisting}

\paragraph{PeTTa accumulator-return compilation.}
A subtlety in PeTTa's compilation of recursive functions: when
a recursive function returns an accumulator variable directly
in one branch, the output can be pinned to the initial value.
The workaround is to wrap the return value in a tagged constructor:

\begin{lstlisting}[language=MeTTa,caption={Streaming loop with tagged return.}]
(= (process-token-stream $state $count)
   (let $result (ds_stream_next)
     (let $tag (ds_stmt_functor $result)
       (case $tag
         ((end
            (let $_ (println! ("Verified" $count "statements."))
              (stream_ok $state)))
          (ns
            (let $stmt (ds_stmt_arg $result 1)
              (pts-continue $stmt $state $count))))))))
\end{lstlisting}

\paragraph{Semantic validation parity.}
The streaming parser (\texttt{next\_statement/2}) performs only
syntactic parsing, while the batch parser
(\texttt{parse\_mm\_file\_compounds/2}) additionally performs 8+
file-wide semantic validations.  To maintain parity,
\texttt{ds\_stream\_init} performs a full batch parse-and-validate
pass first, then tokenizes for streaming.  This two-pass approach
ensures identical accept/reject behavior while still delivering
statements one at a time during verification.

\subsection{pverify.pl: Pure Prolog Verifier}

For comparison, \texttt{pverify.pl} (520 lines) implements the
complete verification algorithm in pure SWI-Prolog, sharing only
\texttt{mm\_primitives.pl} for parsing.  It uses the same state
representation (AVL assoc + ordsets) natively, avoiding all FFI
overhead.

\subsection{pverify\_op\_space: Space-Backed Storage}
\label{sec:opspace}

The space-backed variant replaces \texttt{pverify\_op}'s opaque
Prolog data structures with PeTTa-native specialized spaces,
moving persistent verifier state closer to idiomatic MeTTa.
The design introduces two specialized backends implemented in
\texttt{lib\_zar.pl}:

\begin{itemize}[leftmargin=2em]
  \item \textbf{MapSpace}: key--value storage with AVL sidecar index.
  \item \textbf{SetSpace}: membership storage with ordset sidecar index.
\end{itemize}

Each space provides a dual interface:
(1)~the normal PeTTa space API (\texttt{add-atom},
\texttt{match}, \texttt{get-atoms}, \texttt{remove-atom}) via
Prolog \texttt{multifile} hooks, and
(2)~optimized $O(\log n)$ accessors
(\texttt{sp\_map\_get}, \texttt{sp\_set\_has}, etc.) that read
directly from the sidecar index.
Two construction modes are exposed:
\texttt{canonical} (maintains mirrored atoms) and
\texttt{index\_only} (fast mode, sidecar-only).
\texttt{pverify\_op\_space} uses
\texttt{sp\_new\_map\_fast}/\texttt{sp\_new\_set\_fast} for
high-throughput verification.

The state representation changes from immutable value threading
to mutable space-atom threading at global scope:

\begin{lstlisting}[language=MeTTa,caption={Space-backed state representation.}]
;;; State = (state ConstantsSpace FrameStack LabelsSpace)
;;;   ConstantsSpace: SetSpace atom (mutable)
;;;   FrameStack:     MeTTa list of Frame tuples
;;;   LabelsSpace:    MapSpace atom (mutable)
;;;
;;; Frame = (frame Vars DVs FHyps FLabels EHyps)
;;;   Vars/DVs/FLabels use ordset/assoc hot-path structures
\end{lstlisting}

This hybrid layout is deliberate:
\texttt{Constants} and \texttt{Labels} are persistent global stores
and are excellent space candidates; frame-local fields are short-lived
and remain on ordset/assoc hot paths for speed.
The verifier still walks \texttt{FrameStack} functionally in MeTTa:

\begin{lstlisting}[language=MeTTa,caption={Frame stack traversal with space-backed fields.}]
(= (empty-frame)
   (let $vars (ds_empty_ordset)
     (let $dvs (ds_empty_ordset)
       (let $fl (ds_empty_assoc)
         (frame $vars $dvs () $fl ())))))

(= (lookup-d-key $key $fs)
   (if (== (size-atom $fs) 0)
       False
       (let ($f $rest) (decons-atom $fs)
         (let (frame $v $dvs $fh $fl $eh) $f
           (if (ds_ord_member $key $dvs)
               True
               (lookup-d-key $key $rest))))))
\end{lstlisting}

Crucially, this is still bonafide PeTTa verification:
the semantic core remains in MeTTa functions
(\texttt{process-statement}, \texttt{treat-step},
\texttt{treat-assertion-body}, \texttt{check-dvs},
\texttt{verify-proof}).
Only parser and indexed primitives are delegated to Prolog.
Temporary computation (substitution maps, variable collection,
DV filtering) continues to use \texttt{ds\_*} primitives.

This Phase~1 result establishes \textbf{backend-flexible
verifier state without semantic regression}: 141/141 test
parity, identical architecture for all non-storage code.
Current full-\texttt{set.mm} performance is also near-parity with
\texttt{pverify\_op} (Section~\ref{sec:results}).

\subsection{Why This Is Bonafide MeTTa/PeTTa}
\label{sec:bonafide}

For this project, ``bonafide MeTTa/PeTTa verifier'' means:
\begin{enumerate}[leftmargin=2em]
  \item verifier semantics are in \texttt{.metta},
  \item Prolog is a parser/index backend, not the semantic engine,
  \item storage uses real PeTTa spaces (with normal space interfaces).
\end{enumerate}

The current implementation satisfies all three criteria.
For example, proof-step semantics are dispatched in MeTTa:
\begin{lstlisting}[language=MeTTa,caption={Semantic proof-step dispatch remains in MeTTa.}]
(= (treat-step $entry $stack $state)
   (case $entry
     (((hyp $hkind $stmt) (cons-atom $stmt $stack))
      ((assertion $dvs $fhyps $ehyps $conclusion)
         (treat-assertion-body $dvs $fhyps
           (get-ehyp-stmts-from-pairs $ehyps)
           $conclusion $stack $state)))))
\end{lstlisting}

Likewise, statement handling remains in MeTTa
(\texttt{process-statement}, \texttt{add-c}, \texttt{add-f},
\texttt{add-d}); space operations are invoked as backend primitives,
not as replacements for verifier control logic.

%% ================================================================
\section{Test Suite and Results}
\label{sec:results}

\subsection{Conformance Suite}

The test suite comprises 141 active test cases, plus 10 harness
entries currently skipped because they require features outside
verifier scope:

\begin{itemize}[leftmargin=2em]
  \item \textbf{Positive tests}: Valid \texttt{.mm} files that must
    be accepted (exit code 0)
  \item \textbf{Negative tests}: Invalid files (syntax errors,
    duplicate labels, DV violations, proof failures, etc.) that
    must be rejected (nonzero exit code)
\end{itemize}

Tests are drawn from:
\begin{itemize}[leftmargin=2em]
  \item \texttt{core/small/}: Core Metamath language features
  \item \texttt{unit/}: Focused unit tests for specific edge cases
  \item \texttt{mmverify/}: Tests from the mmverify.py project
  \item \texttt{core/large/}: Large-scale tests including
    \texttt{set.mm}
\end{itemize}

\subsection{Results}

Table~\ref{tab:results} presents the main methodology-chain benchmark.
Each run was executed solo and sequentially on the same machine
(ThinkPad, Linux 6.17, SWI-Prolog 9.x), with explicit memory limits
(\texttt{ulimit -v 6291456}) and one-hour timeout guards.
Prolog baseline used \texttt{swipl --stack\_limit=4g}; PeTTa runs used
\texttt{run.sh} (\texttt{--stack\_limit=8g}).

\begin{table}[ht]
\centering
\caption{Benchmark results on \texttt{set.mm}
  (519K lines, 136{,}289 statements, February 2026).
  Solo sequential runs under fixed memory limits.}
\label{tab:results}
\begin{tabular}{lrrrrr}
\toprule
\textbf{Verifier} & \textbf{Suite} & \textbf{Wall} & \textbf{CPU} & \textbf{ms/stmt} & \textbf{RSS} \\
\midrule
pverify.pl         & 141/141 & 5m 58s  & 343s  & 2.5  & 5.20\,GB \\
pverify\_op        & 141/141 & 28m 32s & 1703s & 12.5 & 3.00\,GB \\
pverify\_op\_space & 141/141 & 28m 48s & 1719s & 12.6 & 3.00\,GB \\
\bottomrule
\end{tabular}
\end{table}
All values in Tables~\ref{tab:results} and~\ref{tab:scaling} are snapshot
measurements from the artifacts listed in the Measurement provenance note
(\S\ref{sec:claims}).

To test scaling before full \texttt{set.mm}, we also measured
\texttt{pverify\_op} and \texttt{pverify\_op\_space} on progressively
larger corpora (Table~\ref{tab:scaling}).

\begin{table}[ht]
\centering
\caption{Intermediate scaling checks for opaque vs space-backed PeTTa.}
\label{tab:scaling}
\begin{tabular}{lrrr}
\toprule
\textbf{Dataset} & \textbf{Lines} & \textbf{pverify\_op (wall)} & \textbf{pverify\_op\_space (wall)} \\
\midrule
\texttt{hol.mm}              & 2{,}251   & 1.35s    & 1.24s \\
\texttt{ql.mm}               & 12{,}533  & 6.55s    & 5.94s \\
\texttt{nf.mm}               & 61{,}205  & 36.50s   & 34.80s \\
\texttt{set.2010-08-29.mm}   & 170{,}876 & 3m 05.52s & 2m 58.87s \\
\texttt{set.mm}              & 519{,}291 & 28m 31.52s & 28m 47.77s \\
\bottomrule
\end{tabular}
\end{table}

On the active small conformance harness, both
\texttt{pverify\_op} and \texttt{pverify\_op\_space} currently report
141/141 with identical 1m\,34s runtime.

%% ================================================================
\section{Performance Analysis}
\label{sec:perf}

For the primary chain (\texttt{pverify.pl} $\rightarrow$
\texttt{pverify\_op} $\rightarrow$ \texttt{pverify\_op\_space}),
the key ratios on \texttt{set.mm} are:
\begin{itemize}[leftmargin=2em]
  \item pure Prolog vs opaque PeTTa:
    $28m31.52s / 5m58.02s \approx 4.78\times$ wall-time,
  \item space-backed vs opaque PeTTa:
    $28m47.77s / 28m31.52s \approx 1.0095\times$ wall-time
    (0.95\% delta).
\end{itemize}
The main result is not that spaces are ``faster'' yet; it is that
state-storage refactoring to spaces can preserve full conformance
and near-identical full-corpus performance.

\subsection{FFI Call Frequency Is the Dominant Cost}

The PeTTa batch verifier performs identical algorithmic work to
the pure Prolog verifier---same AVL tree operations, same ordset
manipulations, same proof stack machine---but every data structure
operation crosses the PeTTa--Prolog FFI boundary.  A single
\texttt{process-statement} call typically invokes 10--30 FFI
functions (\texttt{ds\_assoc\_get}, \texttt{ds\_ord\_add},
\texttt{ds\_stmt\_head}, \texttt{ds\_stmt\_tail}, etc.).
Over 136{,}289 statements, this amounts to millions of FFI
round-trips.

The per-statement overhead quantifies this directly:
2.5\,ms/stmt in pure Prolog vs 12.5\,ms/stmt in PeTTa opaque,
implying $\sim$10\,ms/stmt overhead dominated by FFI granularity.

\subsection{Why Space-Backed Parity Matters}

The space-backed variant demonstrates that specialized spaces can be
used as first-class verifier storage without destabilizing performance
at \texttt{set.mm} scale.
This gives a clean substrate for next-step optimizations (batched
space operations, fewer boundary crossings) while keeping the verifier
logic in MeTTa.

\subsection{Memory Profile}

Interestingly, the PeTTa verifiers use \emph{less} memory
(3.0\,GB) than pure Prolog (5.2\,GB).  In PeTTa, opaque Prolog
terms are referenced indirectly; in pure Prolog, all intermediate
data structures (token lists, parsed statements, verification
state) share a single heap, which grows larger under
SWI-Prolog's non-compacting GC.

\subsection{Benchmark Methodology Note}

An earlier benchmark used concurrent long-running verifier processes,
which produced severely distorted wall times.
All numbers reported here are from solo sequential runs to avoid
contention artifacts.

%% ================================================================
\section{Code Size Comparison}
\label{sec:codesize}

\begin{table}[ht]
\centering
\caption{Lines of code by component.}
\label{tab:loc}
\begin{tabular}{lr}
\toprule
\textbf{Component} & \textbf{Lines} \\
\midrule
mm\_primitives.pl (shared parser)       & 1{,}224 \\
pverify\_ds.pl (FFI data structures)    & 227 \\
lib\_zar.pl (space backends + utils)    & 405 \\
env\_utils.pl (CLI utilities)           & 39 \\
\midrule
pverify.pl (pure Prolog verifier)       & 520 \\
pverify\_op.metta (batch PeTTa)         & 693 \\
pverify\_op\_space.metta (space-backed) & 658 \\
pverify\_op\_stream.metta (streaming)   & 703 \\
\bottomrule
\end{tabular}
\end{table}

The space-backed verifier is \emph{shorter} than the batch version
(658 vs 693 lines) because mutable space operations eliminate the
need to reconstruct state tuples after each mutation---a
\texttt{sp\_map\_put} modifies the space in-place and returns
\texttt{ok}, whereas \texttt{ds\_assoc\_put} returns a new assoc
that must be threaded into a fresh state tuple.  The space backends
themselves (MapSpace + SetSpace + fast/index modes) add substantial
shared infrastructure to
\texttt{lib\_zar.pl},
which is shared library infrastructure reusable by any PeTTa program.

The streaming verifier adds only 10 lines over the batch version---the
\texttt{process-token-stream} loop replaces
\texttt{process-all-statements}, and the FFI import list gains two
entries (\texttt{ds\_stream\_init}, \texttt{ds\_stream\_next}).
The Prolog-side streaming support adds $\sim$30 lines to
\texttt{pverify\_ds.pl}.

%% ================================================================
\section{Lessons Learned}
\label{sec:lessons}

\subsection{PeTTa--Prolog FFI Design}

The opaque data structure pattern---where PeTTa code never inspects
Prolog-side values but manipulates them exclusively through typed
FFI calls---proved highly effective for correctness.  It provides
$O(\log n)$ performance for associative lookups and set operations
while keeping the verification logic in idiomatic MeTTa.
However, current measurements still show about $4.78\times$
for opaque PeTTa vs pure Prolog on \texttt{set.mm}.
This indicates that FFI call frequency, not algorithmic complexity,
is the main performance bottleneck.

\subsection{FFI Boundary Pitfalls}

Two non-obvious FFI issues consumed significant debugging time:
\begin{enumerate}[leftmargin=2em]
  \item \textbf{Parenthesis atoms}: Prolog atoms \texttt{'('}
    and \texttt{')'} are valid Metamath tokens but are interpreted
    as MeTTa grouping when crossing the FFI boundary.  Keeping
    token lists Prolog-resident was the clean solution.
  \item \textbf{Boolean representation}: Prolog's \texttt{true}
    and \texttt{false} atoms do not map to PeTTa's boolean values.
    Using \texttt{case} pattern matching instead of \texttt{if}
    avoids this mismatch.
\end{enumerate}

\subsection{PeTTa Compilation Subtleties}

The accumulator-return pinning issue in recursive functions is a
PeTTa-specific compilation artifact.  The tagged-constructor
workaround (\texttt{stream\_ok \$state}) is simple but non-obvious;
it was diagnosed through systematic binary-search debugging
(progressively simplifying the recursive body until the minimal
failing case was isolated) with assistance from Codex.

\subsection{Human--AI Collaboration}

The pverify family was developed through iterative human--AI
collaboration:
\begin{itemize}[leftmargin=2em]
  \item The human author wrote the original mmverify verifier and
    provided architectural direction, test infrastructure, and
    quality oversight.
  \item AI agents (Claude Code, Codex) performed the bulk of
    implementation, debugging, and optimization work, including
    the opaque data structure refactoring, streaming architecture,
    and the systematic debugging of PeTTa compilation issues.
  \item The 141-test conformance suite served as the shared
    contract between human and AI: no claims of progress without
    passing all tests.
\end{itemize}

%% ================================================================
\section{Future Work: Reducing FFI Overhead}
\label{sec:future}

The performance analysis points to a staged optimization strategy.
Phase~1 is complete and Phase~2 is partially complete.

\paragraph{Phase~1: Specialized spaces (complete).}
\texttt{pverify\_op\_space} replaces opaque Prolog data structures
with PeTTa-native \texttt{MapSpace}/\texttt{SetSpace} backends
(Section~\ref{sec:opspace}).  This achieves backend-flexible
verifier state with 141/141 test parity and no semantic regression
(Table~\ref{tab:results}; provenance note in \S\ref{sec:claims}).
The space backends are generic PeTTa infrastructure
(\texttt{lib\_zar.pl}), not verifier-specific.

\paragraph{Phase~2: Batched hot-path operations.}
This phase is underway: batched substitutions and variable extraction
(\texttt{ds\_apply\_subst}, \texttt{ds\_find\_vars}) are already in
use in the current verifier path.
Next batched targets are substitution-map construction and repeated
label/type lookups in proof loops.

\paragraph{Phase~3: ``Many'' space APIs.}
Adding \texttt{sp\_map\_get\_many}, \texttt{sp\_map\_put\_many},
\texttt{sp\_set\_add\_many} would amortize FFI overhead for bulk
operations like \texttt{foldl-add-c} (adding all constants from
a \texttt{\$c} statement in one call).

\paragraph{Phase~4: Chunked streaming.}
Replace per-statement \texttt{ds\_stream\_next} with a chunked
variant returning $K$ statements per FFI call, reducing streaming
FFI calls from 136{,}289 to $\sim$532 for $K = 256$.

%% ================================================================
\section{Related Work}
\label{sec:related}

\paragraph{Metamath verifiers.}
The canonical reference implementation is \texttt{metamath.exe}
(C, by Megill).  \texttt{mmverify.py} (Python, by Heitzig et al.)
is a concise ($\sim$300 line) verifier widely used as a starting
point.  \texttt{mm0-rs} (Rust, by Carneiro~\cite{carneiro2019mm0})
focuses on performance.  A formally verified Lean~4 verifier
exists~\cite{goertzel2026lean} with a complete
soundness-and-completeness proof.

\paragraph{MeTTa implementations.}
Hyperon Experimental~\cite{goertzel2023metta} is the primary
MeTTa implementation in Rust.  PeTTa compiles MeTTa to SWI-Prolog,
providing mature Prolog libraries through FFI.
Within the implementations and references surveyed for this draft,
pverify appears to be the first non-trivial formal
verification workload implemented in MeTTa.

%% ================================================================
\section{Conclusion}
\label{sec:conclusion}

The pverify family demonstrates that MeTTa, via PeTTa's Prolog
backend, is capable of implementing correct Metamath verification at
full \texttt{set.mm} scale under a controlled methodology chain:
\texttt{mmverify} $\rightarrow$ \texttt{pverify\_hybrid} $\rightarrow$
\texttt{pverify.pl} $\rightarrow$ \texttt{pverify\_op} $\rightarrow$
\texttt{pverify\_op\_space}.
All active implementations keep 141/141 conformance in the recorded
snapshot measurements (\S\ref{sec:claims}).

The core result of this paper is methodological and architectural:
\texttt{pverify\_op\_space} makes verifier storage space-backed and
still stays near performance parity with \texttt{pverify\_op}
(28m\,48s vs 28m\,32s on \texttt{set.mm}, 0.95\% wall delta).
This demonstrates that the verifier can be made more ``space-native''
without sacrificing behavior or practical throughput.

Pure Prolog remains significantly faster (5m\,58s on \texttt{set.mm});
the remaining gap is primarily FFI-boundary granularity rather than
algorithmic mismatch.
That places the next optimization work on batched boundary crossings,
not on rewriting verifier semantics.

%% ================================================================
\begin{thebibliography}{9}

\bibitem{megill2019}
N.~Megill and D.~A.~Wheeler,
\emph{Metamath: A Computer Language for Mathematical Proofs},
Lulu Press, 2019.
Available at \url{http://us.metamath.org/downloads/metamath.pdf}.

\bibitem{carneiro2019mm0}
M.~Carneiro,
``Metamath Zero: Designing a Theorem Prover Prover,''
\emph{arXiv preprint arXiv:1910.10703}, 2019.

\bibitem{goertzel2023metta}
B.~Goertzel \emph{et al.},
``Hyperon: A Framework for AGI at the Human Level and Beyond,''
\emph{arXiv preprint arXiv:2310.18318}, 2023.

\bibitem{petta}
D.~Vos,
``PeTTa: Prolog-based MeTTa Implementation,''
TrueAGI, 2024.

\bibitem{goertzel2026lean}
Z.~Goertzel and Oru\v{z}i,
``Verified Verification: A Lean~4 Proof That the Metamath
Checker Is Correct,'' Working paper, February 2026.

\end{thebibliography}

\end{document}
